\section{问答与对话的综合比较}
\label{sec:nlp-dialogue-comparison}

前五节分别处理了五类研究内容，它们可以共享编码、检索和生成部件，但各自的挑战与完成条件不同。表\ref{tab:nlp-dialogue-lineage}把全章的方法谱系放在同一张表中：每类研究内容按“挑战—方法—解决前者什么问题—代价”展开，同一类内的方法自上而下大致对应演进顺序。读表时可以沿两个方向看：横向看某个方法为解决什么而提出、付出了什么；纵向看同一挑战在问答侧和对话侧的不同处理，例如“证据不足时弃答”与“目标未定时澄清”都是在信息不够时不硬给答案，“引用生成”与“知识支持的回复”都要求输出可追溯到来源。

\begin{BookLongTable}{@{}>{\RaggedRight\arraybackslash}p{17mm}>{\RaggedRight\arraybackslash}p{27mm}>{\RaggedRight\arraybackslash}p{34mm}>{\RaggedRight\arraybackslash}p{\dimexpr.5\linewidth-39mm-8\tabcolsep\relax}>{\RaggedRight\arraybackslash}p{\dimexpr.5\linewidth-39mm\relax}@{}}
\toprule
研究内容 & 挑战 & 方法 & 解决前者的什么问题 & 代价\\
\midrule
\endfirsthead
\multicolumn{5}{@{}l}{\small 表~\thetable（续）}\\
\toprule
研究内容 & 挑战 & 方法 & 解决前者的什么问题 & 代价\\
\midrule
\endhead
\midrule
\multicolumn{5}{r@{}}{\small 续下页}\\
\endfoot
\bottomrule
\caption{问答与对话的方法谱系。每类研究内容内，箭头表示后一方法针对前一方法的缺陷提出；“代价”一列列出新方法引入的主要成本或遗留问题。各方法的细节见第\ref{sec:nlp-dialogue-evidence}至第\ref{sec:nlp-dialogue-collaboration}节。}
\label{tab:nlp-dialogue-lineage}\\
\endlastfoot
    证据问答 & 在给定段落中定位答案并弃答 & BiDAF$\to$\allowbreak BERT跨度阅读 & 预训练联合编码取代手工交互层 & 只在已给窗口内定位\\
     & 大规模文集中找证据 & BM25$\to$\allowbreak DPR$\to$\allowbreak REALM、Atlas & 稠密检索补词项不匹配；联合训练让检索器服务阅读器 & 索引构建与刷新；假负例\\
     & 综合多个片段生成答案 & RAG、FiD & 答案不限于跨度，可融合多片段 & 生成可能无依据\\
     & 查询不准、初检不精、一次检索不够、检索时机固定、答案难追溯 & 查询改写与HyDE、monoBERT重排、IRCoT、Self-RAG、ALCE式带引用生成 & 分别改进检索查询、排序、多跳检索、检索时机与归因 & 调用次数与延迟增加\\
     & 资料能否整体放入上下文 & 长上下文直接阅读、Self-Route & 免去检索遗漏；按需在两条路线间切换 & 位置效应；计算随长度增长\\
     & 证据在表格与数据库中 & SQL与图查询执行、TAPAS、TAGOP、FinQANet & 支持筛选、聚合与多步计算 & 需要模式、算子或程序监督\\
\midrule
    核验 & 主张粒度与证据定位 & FEVER流水线$\to$\allowbreak FActScore分解 & 逐事实判定，避免一处错误拖累整句 & 分解质量决定上限\\
     & 支持、反驳与信息不足 & 可分解注意力、规则与结构化比较、SAFE & 从固定证据扩展到检索证据 & 搜索成本与可复现性\\
     & 据核验结果修订 & RARR、CoVe & 最小编辑保留原回答；自问自答查自相矛盾 & CoVe不引入外部证据\\
\midrule
    任务型对话 & 理解每轮意图与槽位 & CRF槽位标注$\to$\allowbreak JointBERT$\to$\allowbreak 模式约束记录生成 & 联合建模意图与槽位；生成可表达更正 & 生成结果需类型检查\\
     & 跨轮累积与更正需求 & 规则写入$\to$\allowbreak TRADE$\to$\allowbreak SimpleTOD$\to$\allowbreak IC-DST & 复制处理未见值；序列化统一接口；少样本迁移 & 误复制旧值；示例检索敏感\\
     & 选择动作并调用服务 & 有限状态策略$\to$\allowbreak 监督动作预测$\to$\allowbreak 强化学习策略$\to$\allowbreak 大语言模型函数调用 & 从手写规则到从数据和交互中学策略 & 用户模拟器偏差；参数与权限需校验\\
     & 如实告知结果 & 模板回填$\to$\allowbreak SC-LSTM & 语义条件生成减少模板维护 & 软门不保证不遗漏\\
\midrule
    开放对话 & 通用乏味回复 & 检索式$\to$\allowbreak 生成式$\to$\allowbreak MMI重排 & 生成新句；惩罚无条件高概率回复 & 反向模型与候选列表\\
     & 相关性、知识量与技能 & DialoGPT、Meena、BlenderBot & 规模与混合技能数据提高相关性与具体性 & 编造事实、重复与矛盾仍在\\
     & 遵循指令与安全 & 指令微调与偏好对齐的聊天模型 & 按要求回应，减少有害输出 & 标注成本；过度拒答\\
     & 知识、人物与情绪一致 & Transformer Memory Network、Profile Memory Network、用户记忆、MoEL、ESConv & 以知识、人物、情绪或策略为条件 & 条件标注；条件选错向后传播\\
\midrule
    协作对话 & 要求含混 & 规则询问、信息增益、神经EVPI & 从手写假设到学习问题价值 & 只能排序已有候选；用户负担\\
     & 方案难以抽象描述 & 差异候选、成对偏好排序 & 让用户在可见候选中选择并汇总选择 & 阅读负担；独立性假设\\
     & 反馈落到具体位置 & 反馈条件编辑、Self-Refine & 把意见转成局部修改；减少用户逐条反馈 & 改动范围需检查；自评不替代核验\\
\end{BookLongTable}

表中的方法不能用一个统一的“回答质量”比较，因为它们的完成条件不同。比较前要固定问答的资料版本、可见片段与检索和生成预算，固定对话的服务快照、接口权限、用户或模拟器协议、最大轮数与中断条件；允许多种正确表达、等价证据组合或合法交流路径时，在评分前写出等价规则。下面用两个例子说明不同指标的分母怎样决定了结论。

问答侧的分数各自针对不同对象。规范化准确匹配检查答案是否与某个允许参考一致；词元F1按预测与参考的词元多重集合交集计算精确率与召回率。参考为“4月18日／14:00”而预测只有“4月18日”时，精确率1、召回率$1/2$、F1为$2/3$，准确匹配为0。若时间与报名条件都需回答而召回只取得片段A，按两项所需证据计的证据召回率为$1/2$；若生成了两个正确陈述却都引用A，引用支持率仍为$1/2$。答案、证据覆盖和引用三个分数不能互相替代。

对话侧把每轮有效的槽位—值对作为集合，报告槽位F1和联合目标准确率。表\ref{tab:nlp-dialogue-metrics}沿用第\ref{sec:nlp-dialogue-task}节的三轮主例，假设某系统前两轮正确、第三轮错误保留4/17。两边都为空的轮次单独记为正确，非空槽位采用跨轮汇总的微平均。

\begin{table}[htbp]
\centering
\begin{tabular}{cllc}
\toprule
轮次 & 正确状态 & 假设预测状态 & 全状态正确\\
\midrule
1 & 日期、人数均缺失 & 日期、人数均缺失 & 是\\
2 & 4/17，2人 & 4/17，2人 & 是\\
3 & 4/18，2人 & 4/17，2人 & 否\\
\bottomrule
\end{tabular}
\caption{局部状态分数与完整任务的差别。活动编号和身份为固定上下文，不计入可修改槽位；第三轮旧日期使查询继续为空。}
\label{tab:nlp-dialogue-metrics}
\end{table}

三轮共4个正确正值、4个预测正值，匹配3个，微平均精确率、召回率和F1均为$3/4$，联合目标准确率为$2/3$。但第三轮唯一需要变化的日期没有改对，变化项召回率为0；系统可以每轮都给出礼貌回复，却没有查到用户要核对的报名，完整任务成功为0。三个数字都计算无误，回答的却是不同问题：槽位F1问多少槽位值对，联合准确率问多少轮整体对，任务成功问用户的目标是否达成。完整交互还应报告用户补充或更正次数、总轮数、检索与模型调用成本，自反馈调用不记为用户轮次。

比较两个系统时，以相同问题或同一初始对话任务配对。假设四个教学任务上A的最终成功为$(1,1,0,0)$、B为$(1,0,1,0)$，平均成功率都是$1/2$，逐任务差为$(0,1,-1,0)$：两者修复的是不同任务，平均值相同不说明失败机制相同。样本足够时按完整任务配对重采样计算差值区间；多轮对话不能把各轮当成独立样本来虚增样本量。人工选择、Bradley--Terry聚合与第\ref{sec:nlp-eval-judge-protocol}小节的模型评判分别记录，并以独立人工或可执行任务依据校验。诊断时依次检查资料是否取得、本轮是否理解、状态是否更新、动作是否合规、回复是否有据、目标是否达成；给定正确证据或正确状态的诊断实验有助于隔离环节，但不替代端到端比较。
